\begin{proof}[Proof of \cref{obj:lem:hellinger-block-calculus}]
We first establish the finite-reference identity used below. Whenever a real Radon--Nikodym density is used, its value is taken to be zero where the extended-real derivative is infinite; for finite measures, this changes the derivative on a null set.

Let \(\rho_{\mathrm{ref}}\) be a finite measure on the common space \(\Omega\), and suppose measurable nonnegative real densities satisfy
\[
f_{\mathbb B,\mathrm{ref}},f_{\mathbb C,\mathrm{ref}}
   :\Omega\longrightarrow[0,\infty),
\qquad
\mathbb B(dw)=f_{\mathbb B,\mathrm{ref}}(w)\rho_{\mathrm{ref}}(dw),
\qquad
\mathbb C(dw)=f_{\mathbb C,\mathrm{ref}}(w)\rho_{\mathrm{ref}}(dw).
\]
Define
\[
\xi_{\Sigma}:=\mathbb B+\mathbb C,
\qquad
f_{\mathbb B,\Sigma}:=\frac{d\mathbb B}{d\xi_{\Sigma}},
\qquad
f_{\mathbb C,\Sigma}:=\frac{d\mathbb C}{d\xi_{\Sigma}},
\qquad
r_{\Sigma,\mathrm{ref}}:=\frac{d\xi_{\Sigma}}{d\rho_{\mathrm{ref}}},
\]
using the real-density convention above. Radon--Nikodym reconstruction and the chain rule give
\[
\mathbb B(dw)=f_{\mathbb B,\Sigma}(w)\xi_{\Sigma}(dw),
\qquad
\mathbb C(dw)=f_{\mathbb C,\Sigma}(w)\xi_{\Sigma}(dw),
\]
and, \(\rho_{\mathrm{ref}}\)-almost everywhere,
\[
f_{\mathbb B,\mathrm{ref}}
   =f_{\mathbb B,\Sigma}r_{\Sigma,\mathrm{ref}},
\qquad
f_{\mathbb C,\mathrm{ref}}
   =f_{\mathbb C,\Sigma}r_{\Sigma,\mathrm{ref}}.
\]
For these nonnegative densities,
\[
r_{\Sigma,\mathrm{ref}}
 \bigl(\sqrt{f_{\mathbb B,\Sigma}}
             -\sqrt{f_{\mathbb C,\Sigma}}\bigr)^2
=
\bigl(\sqrt{f_{\mathbb B,\Sigma}r_{\Sigma,\mathrm{ref}}}
             -\sqrt{f_{\mathbb C,\Sigma}r_{\Sigma,\mathrm{ref}}}\bigr)^2.
\]
Thus changing measure in the normalization of \cref{obj:def:hellinger} yields
\[
\begin{aligned}
\mathsf h^2(\mathbb B,\mathbb C)
&=\int
 r_{\Sigma,\mathrm{ref}}
 \bigl(\sqrt{f_{\mathbb B,\Sigma}}
             -\sqrt{f_{\mathbb C,\Sigma}}\bigr)^2
 \,d\rho_{\mathrm{ref}}\\
&=\int
 \bigl(\sqrt{f_{\mathbb B,\mathrm{ref}}}
             -\sqrt{f_{\mathbb C,\mathrm{ref}}}\bigr)^2
 \,d\rho_{\mathrm{ref}}.
\end{aligned}
\]
% lean: CausalSmith.Stat.TwosamplePointcateAnnotationFrontier.hellingerSq_real_density

\begin{enumerate}
\item \emph{Finite products.}
Write \(\mathcal U_*\) for the finite indexing set and \(\Omega_{u_*}\) for the measurable space of the pair indexed by \(u_*\). Define
\[
\xi_{u_*}:=\mathbb B_{u_*}+\mathbb C_{u_*},
\qquad
f_{\mathbb B,u_*}:=\frac{d\mathbb B_{u_*}}{d\xi_{u_*}},
\qquad
f_{\mathbb C,u_*}:=\frac{d\mathbb C_{u_*}}{d\xi_{u_*}},
\qquad u_*\in\mathcal U_*,
\]
again as real densities. They are measurable, nonnegative, and integrable, and reconstruction gives
\[
\mathbb B_{u_*}(dw)
   =f_{\mathbb B,u_*}(w)\xi_{u_*}(dw),
\qquad
\mathbb C_{u_*}(dw)
   =f_{\mathbb C,u_*}(w)\xi_{u_*}(dw),
\]
with
\[
\int f_{\mathbb B,u_*}\,d\xi_{u_*}
=
\int f_{\mathbb C,u_*}\,d\xi_{u_*}
=1.
\]
Define the coordinate affinity by
\[
\mathsf a_{\mathrm{aff},u_*}
:=
\int
\sqrt{f_{\mathbb B,u_*}(w)f_{\mathbb C,u_*}(w)}
\,d\xi_{u_*}(w).
\]
The square roots are square-integrable, so their product is integrable. Expanding the square in the finite-reference identity gives
\[
\begin{aligned}
\mathsf h^2(\mathbb B_{u_*},\mathbb C_{u_*})
&=\int
 \left(
 f_{\mathbb B,u_*}+f_{\mathbb C,u_*}
 -2\sqrt{f_{\mathbb B,u_*}f_{\mathbb C,u_*}}
 \right)\,d\xi_{u_*}\\
&=2\bigl(1-\mathsf a_{\mathrm{aff},u_*}\bigr).
\end{aligned}
\]

On the product space, define
\[
\xi_{\otimes}:=\bigotimes_{u_*\in\mathcal U_*}\xi_{u_*},
\qquad
f_{\mathbb B,\otimes}(w)
   :=\prod_{u_*\in\mathcal U_*}f_{\mathbb B,u_*}(w_{u_*}),
\qquad
f_{\mathbb C,\otimes}(w)
   :=\prod_{u_*\in\mathcal U_*}f_{\mathbb C,u_*}(w_{u_*}),
\quad
w\in\prod_{u_*\in\mathcal U_*}\Omega_{u_*}.
\]
To verify the product-density representation, take measurable sets
\[
W_{u_*}\subseteq\Omega_{u_*},
\qquad u_*\in\mathcal U_*.
\]
Product integration gives
\[
\int_{\prod_{u_*}W_{u_*}}
 f_{\mathbb B,\otimes}\,d\xi_{\otimes}
=
\prod_{u_*}
 \int_{W_{u_*}}f_{\mathbb B,u_*}\,d\xi_{u_*}
=
\prod_{u_*}\mathbb B_{u_*}(W_{u_*}).
\]
The same calculation applies to \(\mathbb C\). Uniqueness of finite product measures therefore yields
\[
\left(\bigotimes_{u_*}\mathbb B_{u_*}\right)(dw)
   =f_{\mathbb B,\otimes}(w)\xi_{\otimes}(dw),
\qquad
\left(\bigotimes_{u_*}\mathbb C_{u_*}\right)(dw)
   =f_{\mathbb C,\otimes}(w)\xi_{\otimes}(dw).
\]
In particular,
\[
\int f_{\mathbb B,\otimes}\,d\xi_{\otimes}
=
\int f_{\mathbb C,\otimes}\,d\xi_{\otimes}
=1.
\]
Applying the finite-reference identity and expanding the square as above gives
\[
\mathsf h^2\left(\bigotimes_{u_*}\mathbb B_{u_*},
                \bigotimes_{u_*}\mathbb C_{u_*}\right)
=
2\left(
1-\int\sqrt{f_{\mathbb B,\otimes}f_{\mathbb C,\otimes}}
\,d\xi_{\otimes}
\right).
\]
Nonnegativity of the coordinate densities implies the pointwise identity
\[
\sqrt{f_{\mathbb B,\otimes}(w)f_{\mathbb C,\otimes}(w)}
=
\prod_{u_*}
\sqrt{f_{\mathbb B,u_*}(w_{u_*})
      f_{\mathbb C,u_*}(w_{u_*})}.
\]
Product integration consequently gives
\[
\int\sqrt{f_{\mathbb B,\otimes}f_{\mathbb C,\otimes}}
\,d\xi_{\otimes}
=
\prod_{u_*}\mathsf a_{\mathrm{aff},u_*}.
\]

The affinities satisfy
\[
\mathsf a_{\mathrm{aff},u_*}\ge0,
\qquad
0\le\mathsf h^2(\mathbb B_{u_*},\mathbb C_{u_*})
   =2\bigl(1-\mathsf a_{\mathrm{aff},u_*}\bigr),
\]
and hence
\[
0\le\mathsf a_{\mathrm{aff},u_*}\le1.
\]
For every subset \(\mathcal V_*\subseteq\mathcal U_*\), induction on its cardinality gives
\[
0\le\prod_{u_*\in\mathcal V_*}\mathsf a_{\mathrm{aff},u_*}\le1.
\]
A second induction gives the required affinity-defect bound. The empty product is one and the empty sum is zero. For \(u_*\in\mathcal U_*\setminus\mathcal V_*\), the induction step is
\[
\begin{aligned}
1-\prod_{u'_*\in\mathcal V_*\cup\{u_*\}}
       \mathsf a_{\mathrm{aff},u'_*}
&=
\left(1-\prod_{u'_*\in\mathcal V_*}
       \mathsf a_{\mathrm{aff},u'_*}\right)
+
\bigl(1-\mathsf a_{\mathrm{aff},u_*}\bigr)
 \prod_{u'_*\in\mathcal V_*}\mathsf a_{\mathrm{aff},u'_*}\\
&\le
\sum_{u'_*\in\mathcal V_*}
 \bigl(1-\mathsf a_{\mathrm{aff},u'_*}\bigr)
+
\bigl(1-\mathsf a_{\mathrm{aff},u_*}\bigr).
\end{aligned}
\]
Taking \(\mathcal V_*=\mathcal U_*\) and multiplying by two proves
\[
\begin{aligned}
\mathsf h^2\left(\bigotimes_{u_*}\mathbb B_{u_*},
                \bigotimes_{u_*}\mathbb C_{u_*}\right)
&=2\left(1-\prod_{u_*}\mathsf a_{\mathrm{aff},u_*}\right)\\
&\le2\sum_{u_*}\bigl(1-\mathsf a_{\mathrm{aff},u_*}\bigr)
=\sum_{u_*}\mathsf h^2(\mathbb B_{u_*},\mathbb C_{u_*}).
\end{aligned}
\]
The empty family is included by the stated empty-product and empty-sum conventions.
% lean: CausalSmith.Stat.TwosamplePointcateAnnotationFrontier.hellingerSq_pi_le_sum

\item \emph{A common conditioning marginal.}
Let \(Z\) and \(\Omega\) be the conditioning and outcome spaces. Choose the s-finite reference kernel and jointly measurable densities supplied by the hypothesis, writing them as
\[
\Xi_{\mathrm{ref}}:Z\longrightarrow\{\text{s-finite measures on }\Omega\},
\qquad
f_{\mathbb B,\mathrm{ref}},f_{\mathbb C,\mathrm{ref}}
   :Z\times\Omega\longrightarrow[0,\infty],
\]
so that, for every \(z\in Z\),
\[
\mathbb B_z(dw)
   =f_{\mathbb B,\mathrm{ref}}(z,w)\Xi_{\mathrm{ref},z}(dw),
\qquad
\mathbb C_z(dw)
   =f_{\mathbb C,\mathrm{ref}}(z,w)\Xi_{\mathrm{ref},z}(dw).
\]
Define the sum kernel and normalized densities by
\[
\mathbb K_{\Sigma,z}:=\mathbb B_z+\mathbb C_z
\]
and
\[
\bigl(r_{\mathbb B}(z,w),r_{\mathbb C}(z,w)\bigr)
:=
\begin{cases}
\displaystyle
\left(
\frac{f_{\mathbb B,\mathrm{ref}}(z,w)}
     {f_{\mathbb B,\mathrm{ref}}(z,w)+f_{\mathbb C,\mathrm{ref}}(z,w)},
\frac{f_{\mathbb C,\mathrm{ref}}(z,w)}
     {f_{\mathbb B,\mathrm{ref}}(z,w)+f_{\mathbb C,\mathrm{ref}}(z,w)}
\right),
&
\displaystyle
0<f_{\mathbb B,\mathrm{ref}}(z,w)+f_{\mathbb C,\mathrm{ref}}(z,w)<\infty,
\\[6pt]
(0,0),&\text{otherwise}.
\end{cases}
\]
These are jointly measurable real functions on all of \(Z\times\Omega\), and
\[
0\le r_{\mathbb B}(z,w)\le1,
\qquad
0\le r_{\mathbb C}(z,w)\le1,
\]
because each numerator lies between zero and its positive finite denominator; the other cases have value zero.

For every \(z\),
\[
\mathbb K_{\Sigma,z}(dw)
=
\bigl(f_{\mathbb B,\mathrm{ref}}(z,w)
     +f_{\mathbb C,\mathrm{ref}}(z,w)\bigr)\Xi_{\mathrm{ref},z}(dw),
\]
and the nonnegative extended integral satisfies
\[
\int
\bigl(f_{\mathbb B,\mathrm{ref}}(z,w)
     +f_{\mathbb C,\mathrm{ref}}(z,w)\bigr)
\,d\Xi_{\mathrm{ref},z}(w)
=2.
\]
Thus the sum of the densities is finite \(\Xi_{\mathrm{ref},z}\)-almost everywhere. Where that sum is zero, both densities are zero. Cancellation therefore gives, \(\Xi_{\mathrm{ref},z}\)-almost everywhere,
\[
\begin{aligned}
\bigl(f_{\mathbb B,\mathrm{ref}}+f_{\mathbb C,\mathrm{ref}}\bigr)(z,w)
\,r_{\mathbb B}(z,w)
&=f_{\mathbb B,\mathrm{ref}}(z,w),\\
\bigl(f_{\mathbb B,\mathrm{ref}}+f_{\mathbb C,\mathrm{ref}}\bigr)(z,w)
\,r_{\mathbb C}(z,w)
&=f_{\mathbb C,\mathrm{ref}}(z,w).
\end{aligned}
\]
Consequently, for every \(z\),
\[
\mathbb B_z(dw)=r_{\mathbb B}(z,w)\mathbb K_{\Sigma,z}(dw),
\qquad
\mathbb C_z(dw)=r_{\mathbb C}(z,w)\mathbb K_{\Sigma,z}(dw).
\]

Define the finite joint reference measure by
\[
\Xi_{\Sigma}(dz,dw):=\nu(dz)\mathbb K_{\Sigma,z}(dw).
\]
The kernel density representations and joint measurability imply
\[
\nu(dz)\mathbb B_z(dw)
   =r_{\mathbb B}(z,w)\Xi_{\Sigma}(dz,dw),
\qquad
\nu(dz)\mathbb C_z(dw)
   =r_{\mathbb C}(z,w)\Xi_{\Sigma}(dz,dw).
\]
Moreover,
\[
\Xi_{\Sigma}(Z\times\Omega)
=\int\mathbb K_{\Sigma,z}(\Omega)\,d\nu(z)=2.
\]
The normalized-density bounds give
\[
0\le\sqrt{r_{\mathbb B}(z,w)}\le1,
\qquad
0\le\sqrt{r_{\mathbb C}(z,w)}\le1,
\]
and hence
\[
0\le
\bigl(\sqrt{r_{\mathbb B}(z,w)}
      -\sqrt{r_{\mathbb C}(z,w)}\bigr)^2
\le4.
\]
The measurable square-root difference is therefore integrable against \(\Xi_{\Sigma}\). Applying the finite-reference identity and then integration over the kernel gives
\[
\begin{aligned}
&\mathsf h^2\bigl(\nu(dz)\mathbb B_z(dw),
                 \nu(dz)\mathbb C_z(dw)\bigr)\\
&\quad=
\int
\bigl(\sqrt{r_{\mathbb B}(z,w)}
      -\sqrt{r_{\mathbb C}(z,w)}\bigr)^2
\,d\Xi_{\Sigma}(z,w)\\
&\quad=
\int_Z
\left[
\int_\Omega
\bigl(\sqrt{r_{\mathbb B}(z,w)}
      -\sqrt{r_{\mathbb C}(z,w)}\bigr)^2
\,d\mathbb K_{\Sigma,z}(w)
\right]\,d\nu(z).
\end{aligned}
\]
For each \(z\), the same finite-reference identity and the conditional density representations yield
\[
\int_\Omega
\bigl(\sqrt{r_{\mathbb B}(z,w)}
      -\sqrt{r_{\mathbb C}(z,w)}\bigr)^2
\,d\mathbb K_{\Sigma,z}(w)
=
\mathsf h^2(\mathbb B_z,\mathbb C_z).
\]
Substitution proves
\[
\mathsf h^2\bigl(\nu(dz)\mathbb B_z(dw),
                \nu(dz)\mathbb C_z(dw)\bigr)
=
\int\mathsf h^2(\mathbb B_z,\mathbb C_z)\,d\nu(z).
\]
% lean: CausalSmith.Stat.TwosamplePointcateAnnotationFrontier.hellingerSq_compProd_common

\item \emph{A common independent probability factor.}
For the given \(\mathbb B,\mathbb C\), use
\[
\xi_{\Sigma}:=\mathbb B+\mathbb C,
\qquad
f_{\mathbb B,\Sigma}:=\frac{d\mathbb B}{d\xi_{\Sigma}},
\qquad
f_{\mathbb C,\Sigma}:=\frac{d\mathbb C}{d\xi_{\Sigma}},
\]
with the real-density convention above. The sum of the product laws is
\[
(\mathbb B\otimes\nu)+(\mathbb C\otimes\nu)
=\xi_{\Sigma}\otimes\nu.
\]
The product Radon--Nikodym identities give, \((\xi_{\Sigma}\otimes\nu)\)-almost everywhere,
\[
\frac{d(\mathbb B\otimes\nu)}{d(\xi_{\Sigma}\otimes\nu)}(w,z)
=f_{\mathbb B,\Sigma}(w),
\qquad
\frac{d(\mathbb C\otimes\nu)}{d(\xi_{\Sigma}\otimes\nu)}(w,z)
=f_{\mathbb C,\Sigma}(w).
\]
Substituting into the defining square-root formula and integrating a function of the first coordinate gives
\[
\begin{aligned}
\mathsf h^2(\mathbb B\otimes\nu,\mathbb C\otimes\nu)
&=
\int_{\Omega\times Z}
\bigl(\sqrt{f_{\mathbb B,\Sigma}(w)}
      -\sqrt{f_{\mathbb C,\Sigma}(w)}\bigr)^2
\,d(\xi_{\Sigma}\otimes\nu)(w,z)\\
&=
\nu(Z)
\int_\Omega
\bigl(\sqrt{f_{\mathbb B,\Sigma}(w)}
      -\sqrt{f_{\mathbb C,\Sigma}(w)}\bigr)^2
\,d\xi_{\Sigma}(w)\\
&=\mathsf h^2(\mathbb B,\mathbb C),
\end{aligned}
\]
where the last equality uses \(\nu(Z)=1\).
% lean: CausalSmith.Stat.TwosamplePointcateAnnotationFrontier.hellingerSq_prod_probability
\end{enumerate}
\end{proof}
